\documentclass[UTF8,a4paper,12pt,oneside,openany,fontset=fandol]{ctexbook}
\usepackage{amsmath,amssymb,amsthm,bm,esint,graphicx,geometry,fancyhdr,enumitem,array}
\usepackage{tikz}
\usetikzlibrary{arrows.meta,calc,angles,quotes}
\tikzset{
  vec/.style={thick,-{Stealth[length=5pt,width=3.5pt]}},
  dim/.style={very thin,{Stealth[length=3.5pt,width=2.5pt]}-{Stealth[length=3.5pt,width=2.5pt]}},
  dashedaux/.style={thin,dashed}
}
\usepackage[unicode,hidelinks]{hyperref}

\DeclareMathOperator{\tr}{tr}
\DeclareMathOperator{\erf}{erf}
\DeclareMathOperator{\sh}{sh}
\DeclareMathOperator{\ch}{ch}
\DeclareMathOperator{\tg}{tg}

\geometry{left=22mm,right=22mm,top=20mm,bottom=20mm,headheight=16pt}
\pagestyle{fancy}
\fancyhf{}
\fancyhead[L]{\small 杨孔庆数学物理方法做题本}
\fancyhead[R]{\small\nouppercase{\leftmark}}
\fancyfoot[C]{\thepage}
\renewcommand{\chaptermark}[1]{\markboth{第\thechapter 章\quad #1}{}}

\setlength{\parindent}{2em}
\setlength{\emergencystretch}{3em}
\linespread{1.15}
\ctexset{
  part={name={第,篇},number=\chinese{part}},
  chapter={beforeskip=8pt,afterskip=16pt}
}
\setlist[enumerate]{label=（\arabic*）,leftmargin=2.5em,itemsep=2pt,topsep=3pt}

\newcommand{\problem}[2]{%
  \par\phantomsection\pdfbookmark[2]{题 #1}{problem-#1}%
  \noindent\textbf{#1}\quad #2\par
}
\newcommand{\sproblem}[2]{%
  \par\phantomsection\pdfbookmark[2]{*题 #1}{problem-#1}%
  \noindent\textbf{\textsuperscript{*}#1}\quad #2\par
}

\newsavebox{\figbox}
\newlength{\maxfigheight}
\newcommand{\drawfig}[2]{%
  \begingroup
  \setlength{\maxfigheight}{42mm}%
  \sbox{\figbox}{\resizebox{#1\linewidth}{!}{\input{tikzfigs/fig_#2.tex}}}%
  \ifdim\dimexpr\ht\figbox+\dp\figbox\relax>\maxfigheight
    \resizebox{!}{\maxfigheight}{\usebox{\figbox}}%
  \else
    \usebox{\figbox}%
  \fi
  \endgroup
}
\def\figlabel#1_#2\relax{#1.#2}
\newcommand{\fig}[2][0.48]{%
  \par\addvspace{4pt}%
  {\centering\drawfig{#1}{#2}\par\vspace{2pt}{\footnotesize 习题\ \figlabel#2\relax\ 图}\par}%
  \addvspace{2pt}%
}
\newcommand{\problemfig}[4][0.36]{%
  \par\phantomsection\pdfbookmark[2]{题 #2}{problem-#2}%
  \noindent\begin{minipage}[t]{\dimexpr\linewidth-#1\linewidth-6mm\relax}%
    \vspace{0pt}\textbf{#2}\quad #3%
  \end{minipage}\hfill%
  \begin{minipage}[t]{#1\linewidth}%
    \vspace{0pt}\centering\drawfig{1}{#4}\par\vspace{2pt}{\footnotesize 习题\ \figlabel#4\relax\ 图}%
  \end{minipage}\par
}

\newcommand{\answerarea}{\par\vfill}

\begin{document}
\begin{titlepage}
  \centering
  \vspace*{.24\textheight}
  {\Huge\bfseries 杨孔庆《数学物理方法》\\[0.6em]做题本\par}
  \vspace{1.8cm}
  {\Large 习题整理与作答空间\par}
  \vspace{0.8cm}
  {\large 第一至十一、十三至十九、二十二章\quad 每页两题\par}
  \vfill
  {\large 姓名：\underline{\hspace{4cm}}\quad 日期：\underline{\hspace{3cm}}\par}
\end{titlepage}

\frontmatter
\tableofcontents
\clearpage

\mainmatter
\part{线性空间及线性算子}
\chapter{$\mathbf{R}^3$ 空间的向量分析}

\problem{1.1}{$\boldsymbol{C}$ 为 $\mathbf{R}^3$ 空间的已知向量，若有向量 $\boldsymbol{A}$ 和 $\boldsymbol{B}$，且 $\boldsymbol{A}+\boldsymbol{B}=\boldsymbol{C}$ 和 $\boldsymbol{A}-\boldsymbol{B}=\boldsymbol{C}$，求 $\boldsymbol{A},\boldsymbol{B}$ 向量。}
\answerarea

\problem{1.2}{$\boldsymbol{C}$ 和 $\boldsymbol{D}$ 为 $\mathbf{R}^3$ 空间的已知向量，若有向量 $\boldsymbol{A}$ 和 $\boldsymbol{B}$，且 $\boldsymbol{A}+\boldsymbol{B}=\boldsymbol{C}$ 和 $\boldsymbol{A}-\boldsymbol{B}=\boldsymbol{D}$，求 $\boldsymbol{A},\boldsymbol{B}$ 向量。

如果在笛卡儿坐标系中 $\boldsymbol{C}=c_i\boldsymbol{e}_i,\,\boldsymbol{D}=d_i\boldsymbol{e}_i\ (i=1,2,3)$，请写出 $\boldsymbol{A},\boldsymbol{B}$ 向量的具体表达式。}
\answerarea
\clearpage

\problemfig[0.36]{1.3}{有向量 $\boldsymbol{r}_1,\boldsymbol{r}_2$ 和 $\boldsymbol{r}_3$ 如习题 1.3 图所示，请证明余弦定理
\[
  r_3^2 = r_1^2 + r_2^2 - 2r_1r_2\cos\theta
\]
其中 $|\boldsymbol{r}_1|=r_1,\,|\boldsymbol{r}_2|=r_2,\,|\boldsymbol{r}_3|=r_3$。}{1_3}
\answerarea

\problem{1.4}{请用向量运算（不用坐标系）直接证明
\[
  (\boldsymbol{A}+\boldsymbol{B})\cdot(\boldsymbol{A}-\boldsymbol{B}) = A^2 - B^2
\]
请用 Cartesian 坐标系再证明一遍上述结论。}
\answerarea
\clearpage

\problem{1.5}{在 $\mathbf{R}^3$ 空间中，在空间转动 $\boldsymbol{A}$ 的变换下（即直角坐标系的转动变换），向量 $\boldsymbol{X}$ 变为 $\boldsymbol{X}'$：
\[
  \boldsymbol{X}' = \boldsymbol{A}\boldsymbol{X}
\]
且满足 $\boldsymbol{A}\cdot\boldsymbol{A}^{\mathrm{T}}=\boldsymbol{I}$，$\boldsymbol{A}^{\mathrm{T}}$ 为 $\boldsymbol{A}$ 的转置，$\boldsymbol{I}$ 为单位矩阵，请证明：两个向量 $\boldsymbol{X}$ 和 $\boldsymbol{Y}$ 的标量积 $\boldsymbol{X}\cdot\boldsymbol{Y}$ 在此转动变换下为不变量。}
\answerarea

\problem{1.6}{请证明，在 $\mathbf{R}^3$ 空间中向量 $\boldsymbol{A}$ 和 $\boldsymbol{B}$ 满足
\[
  (\boldsymbol{A}\times\boldsymbol{B})\cdot(\boldsymbol{A}\times\boldsymbol{B}) = A^2B^2 - (\boldsymbol{A}\cdot\boldsymbol{B})^2
\]}
\answerarea
\clearpage

\problem{1.7}{在 $\mathbf{R}^3$ 空间中证明
\begin{enumerate}
  \item $\nabla\cdot\boldsymbol{r}=3$；
  \item $\nabla\times\boldsymbol{r}=0$；
  \item $\nabla\cdot(\boldsymbol{A}+\boldsymbol{B})=\nabla\cdot\boldsymbol{A}+\nabla\cdot\boldsymbol{B}$；
  \item $\nabla\times(\boldsymbol{A}+\boldsymbol{B})=\nabla\times\boldsymbol{A}+\nabla\times\boldsymbol{B}$。
\end{enumerate}}
\answerarea

\problem{1.8}{请证明：静电场是一个无旋场，静磁场是一个无源（点源）场。}
\answerarea
\clearpage

\problem{1.9}{若 $\varphi(\boldsymbol{r})$ 为 $\mathbf{R}^3$ 空间中的标量函数，$\boldsymbol{A}(\boldsymbol{r})$ 和 $\boldsymbol{B}(\boldsymbol{r})$ 为此空间中的向量函数，请证明：
\begin{enumerate}
  \item $\nabla\cdot(\varphi\boldsymbol{A})=\boldsymbol{A}\cdot(\nabla\varphi)+\varphi\nabla\cdot\boldsymbol{A}$；
  \item $\nabla\times(\varphi\boldsymbol{A})=\varphi\nabla\times\boldsymbol{A}-\boldsymbol{A}\times(\nabla\varphi)$；
  \item $\nabla\times(\boldsymbol{A}\times\boldsymbol{B})=(\boldsymbol{B}\cdot\nabla)\boldsymbol{A}-\boldsymbol{B}(\nabla\cdot\boldsymbol{A})-(\boldsymbol{A}\cdot\nabla)\boldsymbol{B}+\boldsymbol{A}(\nabla\cdot\boldsymbol{B})$；
  \item $\nabla(\boldsymbol{A}\cdot\boldsymbol{B})=(\boldsymbol{B}\cdot\nabla)\boldsymbol{A}+(\boldsymbol{A}\cdot\nabla)\boldsymbol{B}+\boldsymbol{B}\times\nabla\times\boldsymbol{A}+\boldsymbol{A}\times\nabla\times\boldsymbol{B}$。
\end{enumerate}}
\answerarea

\problem{1.10}{请计算：
\[
  \nabla^2\frac{1}{r}
\]
其中 $r=|\boldsymbol{r}|$。}
\answerarea
\clearpage

\problem{1.11}{在 $\mathbf{R}^3$ 空间中，对于定义在单连通区域 $\Omega$ 上的向量场 $\boldsymbol{A}(\boldsymbol{r})$，若 $\boldsymbol{A}$ 在 $\Omega$ 的边界 $\partial\Omega$ 上的法向值完全确定，且 $\boldsymbol{A}$ 在 $\Omega$ 上的散度和旋度存在：
\[
  \nabla\cdot\boldsymbol{A}=q,\qquad \nabla\times\boldsymbol{A}=\boldsymbol{H}
\]
请证明，在 $\Omega$ 上这样的向量场 $\boldsymbol{A}$ 是唯一的。这一结论被称为 Helmholtz 定理。（提示：设另一个与 $\boldsymbol{A}$ 完全满足相同条件的向量场 $\boldsymbol{B}$，证明 $\boldsymbol{B}=\boldsymbol{A}$。）}
\answerarea

\problem{1.12}{若刚体以恒定角速度 $\boldsymbol{\omega}$ 转动，有一粒子固定于刚体以速度 $\boldsymbol{v}$ 运动，求证
\[
  \nabla\times\boldsymbol{v}=2\boldsymbol{\omega}
\]
（提示：由 $\boldsymbol{v}=\boldsymbol{v}_C+\boldsymbol{\omega}\times\boldsymbol{r}$，取质心 $C$ 为参考点，建立坐标系，在此坐标系中 $\boldsymbol{v}_C=0$，进行运算。）}
\answerarea
\clearpage

\problem{1.13}{由 Ampere（安培）定律
\[
  \int_S \boldsymbol{J}\cdot\mathrm{d}\boldsymbol{\sigma} = \int_{\partial S}\boldsymbol{H}\cdot\mathrm{d}\boldsymbol{l}
\]
其中 $\boldsymbol{J}$ 为穿过空间曲面 $S$ 的电流密度向量，$\mathrm{d}\boldsymbol{\sigma}$ 为 $S$ 的有向面元，其法向正向由右手定则确定，$\boldsymbol{H}$ 为磁场强度向量，$\mathrm{d}\boldsymbol{l}$ 为 $S$ 的边界 $\partial S$ 的线元，其方向由右手定则确定，它与 $\mathrm{d}\boldsymbol{\sigma}$ 的定向自洽。

请证明：$\nabla\times\boldsymbol{H}=\boldsymbol{J}$。（提示：应用 Stokes 公式证明。）}
\answerarea

\problem{1.14}{有一物理体系，它由三个向量场 $\boldsymbol{E},\boldsymbol{H},\boldsymbol{A}$ 和一个实标量场 $V$ 来描述，并且 $\boldsymbol{E},\boldsymbol{H},\boldsymbol{A}$ 和 $V$ 满足如下关系：
\[
  \nabla\cdot\boldsymbol{E}=-\mu^2 V,\quad \boldsymbol{E}=\frac{\partial\boldsymbol{A}}{\partial t}-\nabla V,\quad \boldsymbol{H}=\nabla\times\boldsymbol{A},\quad \nabla\times\boldsymbol{H}=\frac{\partial\boldsymbol{E}}{\partial t}-\mu^2\boldsymbol{A}
\]
其中 $\mu$ 为正的常数。求证：
\begin{enumerate}
  \item $\nabla\cdot\boldsymbol{A}+\dfrac{\partial V}{\partial t}=0$；
  \item $\nabla^2 V-\mu^2 V+\dfrac{\partial^2 V}{\partial t^2}=0$；
  \item 当我们考虑在区域 $\Omega$ 中，$V$ 与时间无关，且在 $\Omega$ 的边界 $\partial\Omega$ 上有 $V\big|_{\partial\Omega}=0$ 时，请证明：在 $\Omega$ 中 $V=0$。
\end{enumerate}}
\answerarea
\clearpage

\problem{1.15}{若定义算符 $\hat{\boldsymbol{p}}=-\mathrm{i}\hbar\nabla$，$\hat{\boldsymbol{B}}=\nabla\times\boldsymbol{A}$，其中 $\mathrm{i}$ 为虚数单位，$\hbar$ 为常数，$\boldsymbol{A}$ 为向量场。请证明：
\[
  (\hat{\boldsymbol{p}}-e\boldsymbol{A})\times(\hat{\boldsymbol{p}}-e\boldsymbol{A})=\mathrm{i}e\hbar\hat{\boldsymbol{B}}
\]
其中 $e$ 为常数。}
\answerarea
\clearpage

\chapter{$\mathbf{R}^3$ 空间曲线坐标系中的向量分析}

\problem{2.1}{在柱坐标 $(\rho,\varphi,z)$ 下，试计算其坐标基 $\boldsymbol{e}_\rho,\boldsymbol{e}_\varphi,\boldsymbol{e}_z$（可称为空间中的标架场）的散度和旋度，即
\[
  \nabla\cdot\boldsymbol{e}_\rho,\quad \nabla\cdot\boldsymbol{e}_\varphi,\quad \nabla\cdot\boldsymbol{e}_z \quad\text{和}\quad \nabla\times\boldsymbol{e}_\rho,\quad \nabla\times\boldsymbol{e}_\varphi,\quad \nabla\times\boldsymbol{e}_z
\]
所得的计算结果，能否说明它们的物理意义？}
\answerarea

\problem{2.2}{在球坐标 $(r,\theta,\varphi)$ 下，试计算其坐标基矢 $\boldsymbol{e}_r,\boldsymbol{e}_\theta,\boldsymbol{e}_\varphi$ 所构成的标架场的散度和旋度，即
\[
  \nabla\cdot\boldsymbol{e}_r,\quad \nabla\cdot\boldsymbol{e}_\theta,\quad \nabla\cdot\boldsymbol{e}_\varphi \quad\text{和}\quad \nabla\times\boldsymbol{e}_r,\quad \nabla\times\boldsymbol{e}_\theta,\quad \nabla\times\boldsymbol{e}_\varphi
\]
所得的计算结果能否用物理学中的某些场量给予恰当的解释？}
\answerarea
\clearpage

\problem{2.3}{若有一个力场在直角坐标系下的表达式为
\[
  \boldsymbol{F} = -\frac{y}{x^2+y^2}\boldsymbol{e}_x + \frac{x}{x^2+y^2}\boldsymbol{e}_y
\]
\begin{enumerate}
  \item 请用柱坐标系把此力场表示出来；
  \item 在柱坐标系下计算此力场的散度和旋度；
  \item 若有一物体在此力场中沿 $\boldsymbol{e}_\rho$ 方向运动一单位长度时，求此力场 $\boldsymbol{F}$ 所做的功；
  \item 若有一物体在此力场的作用下，沿 $x$-$y$ 平面的一单位圆顺时针运动一圈，求此力场 $\boldsymbol{F}$ 所做的功，并分析此功与力场 $\boldsymbol{F}$ 的旋度的关系。
\end{enumerate}}
\answerarea

\problem{2.4}{请证明
\begin{enumerate}
  \item $\nabla r^n = n r^{n-2}\boldsymbol{r}$；
  \item $\nabla^2 r^n = n(n+1) r^{n-2}$，
\end{enumerate}
其中 $r=|\boldsymbol{r}|$。}
\answerarea
\clearpage

\problem{2.5}{假设点电荷 $q$ 的势为
\[
  \psi(\boldsymbol{r}) = \frac{q}{r^{1+\varepsilon}},\quad 0<\varepsilon<1
\]
在球坐标系下，计算此势（在 $r\neq 0$ 处）的场强 $\boldsymbol{E}\ (\boldsymbol{E}=-\nabla\psi)$ 并计算：
\begin{enumerate}
  \item $\nabla\cdot\boldsymbol{E}$；
  \item $\nabla\times\boldsymbol{E}$；
  \item 对半径为 $R$ 的球形区域，推导类似 Gauss 定律的结果，对于此结果，检验当 $\varepsilon\to 0$ 时的极限情况，并同点电荷的 Coulomb 势 $\psi(\boldsymbol{r})=\dfrac{q}{r}$ 做比较，说明这一电势的改变所引起的物理结果的变化。
\end{enumerate}}
\answerarea

\problem{2.6}{在球坐标系下，已知向量场 $\boldsymbol{A}(\boldsymbol{r})$ 的旋度为
\[
  \nabla\times\boldsymbol{A} = \frac{1}{r^2}\boldsymbol{e}_r
\]
求此向量场 $\boldsymbol{A}$。所得的场 $\boldsymbol{A}$ 是唯一的吗？如果不唯一，请说明理由。}
\answerarea
\clearpage

\chapter{线性空间}

\problem{3.1}{设 $A,B$ 都为 $n$ 维空间中的坐标变换的矩阵，即都为 $n\times n$ 的非奇异（$\det A\neq 0,\ \det B\neq 0$）方阵，求证
\[
  \det(A\cdot B) = \det A \det B
\]}
\answerarea

\problem{3.2}{试证明，若 $A$ 为实内积空间的一个自伴算子，则当且仅当对此空间内所有的向量 $X$，有 $\langle X,AX\rangle=0$ 时，$A=0$，其中 $\langle\ ,\ \rangle$ 为内积。}
\answerarea
\clearpage

\problem{3.3}{若 $A$ 为 $n$ 维空间中的反对称算子（即 $n\times n$ 的反对称矩阵），请证明
\[
  \tr A = 0
\]
（$\tr A$ 为矩阵 $A$ 的迹。）}
\answerarea

\problem{3.4}{若算子 $A$ 为 $3\times 3$ 的矩阵
\[
  A = \begin{pmatrix} 1 & 2 & 3 \\ 2 & 2 & 1 \\ 3 & 4 & 3 \end{pmatrix}
\]
求 $A$ 的伴随算子 $A^+$ 和逆算子 $A^{-1}$。}
\answerarea
\clearpage

\problem{3.5}{求下列矩阵的本征值及相应的本征向量
\begin{enumerate}
  \item $A=\begin{pmatrix} 0 & 1 & 0 \\ 1 & 0 & 0 \\ 0 & 0 & 0 \end{pmatrix}$ \qquad
  \item $A=\begin{pmatrix} 1 & 0 & 0 \\ 0 & 0 & 1 \\ 0 & 1 & 0 \end{pmatrix}$
  \item $A=\begin{pmatrix} 2 & 0 & 0 \\ 0 & 1 & 1 \\ 0 & 1 & 1 \end{pmatrix}$ \qquad
  \item $A=\begin{pmatrix} 1 & 1 & 1 \\ 1 & 1 & 1 \\ 1 & 1 & 1 \end{pmatrix}$
\end{enumerate}}
\answerarea

\problem{3.6}{已知矩阵 $A=\begin{pmatrix} -1 & -2 & 2 \\ 0 & 1 & 0 \\ 0 & 0 & 1 \end{pmatrix}$，求 $A$ 的本征值和相应的本征向量。若此方阵可对角化，求出使 $A$ 对角化的相似变换 $MAM^{-1}$ 的满秩矩阵 $M$。}
\answerarea
\clearpage

\problem{3.7}{求 Pauli（泡利）矩阵
\[
  \sigma_x = \begin{pmatrix} 0 & 1 \\ 1 & 0 \end{pmatrix},\quad
  \sigma_y = \begin{pmatrix} 0 & -\mathrm{i} \\ \mathrm{i} & 0 \end{pmatrix},\quad
  \sigma_z = \begin{pmatrix} 1 & 0 \\ 0 & -1 \end{pmatrix}
\]
的本征值及相应的本征向量，其中 $\mathrm{i}$ 为虚数单位。}
\answerarea
\clearpage


\part{复变函数}
\chapter{复变函数的概念}

\problem{4.1}{写出下列复数的实部、虚部、模和幅角。
\begin{enumerate}
  \item $1-\mathrm{i}\sqrt{3}$ \qquad
  \item $\ln\mathrm{i}$ \qquad
  \item $(1+\mathrm{i})^{1-\mathrm{i}}$ \qquad
  \item $\sqrt[\mathrm{i}]{\mathrm{i}}$
  \item $\sin(a+\mathrm{i}b),\quad a,b\in R$ \qquad
  \item $\sqrt{a+\mathrm{i}b},\quad a,b\in R$
\end{enumerate}}
\answerarea

\problem{4.2}{写出下列复变函数 $W=f(z)=u(x,y)+\mathrm{i}v(x,y)$ 的实部 $u(x,y)$、虚部 $v(x,y)$、模 $|f(z)|$ 和幅角 $\arg f(z)$。
\begin{enumerate}
  \item $\mathrm{e}^z$ \qquad
  \item $\mathrm{e}^{\mathrm{i}z}$ \qquad
  \item $\mathrm{e}^{\mathrm{i}\sin x},\quad x\in R$
  \item $z^3$ \qquad
  \item $\sqrt{\mathrm{i}+\mathrm{i}2x\sqrt{x^2-1}},\quad x\in R$
\end{enumerate}}
\answerarea
\clearpage

\problem{4.3}{求下列方程的解，并把解在复平面上用图形表示出来，说明它们的几何意义。
\begin{enumerate}
  \item $|z|=1$ \qquad
  \item $|z|<1$
  \item $|z|\geqslant 1$ \qquad
  \item $1<|z|<2$
  \item $|z+\mathrm{i}|\leqslant |z-\mathrm{i}|$ \qquad
  \item $|z-a|=|z-b|\quad a,b\text{ 为复平面上的点}$
  \item $0<\arg(z-1)<\dfrac{\pi}{4}$ \qquad
  \item $\operatorname{Re}\!\left(\dfrac{1}{z}\right)=\dfrac{1}{2}$
  \item $|z_1+z_2|^2+|z_1-z_2|^2=2|z_1|^2+2|z_2|^2$
  \item $\sin z=0$ \qquad
  \item $\cos z=0$
  \item $\sh z=0$ \qquad
  \item $\ch z=0$
\end{enumerate}}
\answerarea

\problem{4.4}{试证明：
\begin{enumerate}
  \item $|\sin z|\geqslant \sin x$ \qquad
  \item $|\cos z|\geqslant |\cos x|$
\end{enumerate}}
\answerarea
\clearpage

\problem{4.5}{对于在实变量 $x$ 中成立的关系式
\[
  \sin(x_1+x_2) = \sin x_1\cos x_2 + \sin x_2\cos x_1
\]
请证明这一关系式在复变量的情况下也成立，即
\[
  \sin(z_1+z_2) = \sin z_1\cos z_2 + \sin z_2\cos z_1
\]}
\answerarea

\problem{4.6}{对下列复变函数作出 Riemann 面的示意图，并指出函数的枝点及其阶数，同时说明你所作的 Riemann 面的割线的取法。
\begin{enumerate}
  \item $\sqrt{(z-a)(z-b)}$ \qquad
  \item $\sqrt[3]{(z-a)(z-b)}$ \qquad
  \item $\ln(z-a)$
  \item $\sqrt{(z-a)(z-b)(z-c)}$
\end{enumerate}
其中 $a,b,c$ 为互不相等的复常数。}
\answerarea
\clearpage

\chapter{解析函数}

\problem{5.1}{请说明下面复变函数在复平面上的哪些点或区域上可导，在复平面上的哪些区域上解析。
\begin{enumerate}
  \item $f(z)=\operatorname{Re} z$ \qquad
  \item $f(z)=z^*$ \qquad
  \item $f(z)=|z|$
  \item $f(z)=\mathrm{e}^z$ \qquad
  \item $f(z)=z^2$
  \item $f(z)=(x^2-y^2-x)+\mathrm{i}(2xy-y^2)$ \qquad
  \item $f(z)=x^3-y^3-\mathrm{i}2x^2y$
\end{enumerate}}
\answerarea

\problem{5.2}{若 $f(z)=u(x,y)+\mathrm{i}v(x,y)$ 和 $f^*(z)=u(x,y)-\mathrm{i}v(x,y)$ 都为解析函数，请说明 $f(z)$ 的函数形式。}
\answerarea
\clearpage

\problem{5.3}{若 $f(z)=u(x,y)+\mathrm{i}v(x,y)$ 为解析函数，请证明 $f^*(z^*)$ 也是解析函数。}
\answerarea

\problem{5.4}{请证明：柱坐标系下的解析函数
\[
  f(z) = u(\rho,\varphi) + \mathrm{i}v(\rho,\varphi)
\]
满足的 Cauchy-Riemann 方程
\[
  \left\{
  \begin{aligned}
    \frac{\partial u}{\partial \rho} &= \frac{1}{\rho}\frac{\partial v}{\partial \varphi} \\
    \frac{1}{\rho}\frac{\partial u}{\partial \varphi} &= -\frac{\partial v}{\partial \rho}
  \end{aligned}
  \right.
\]}
\answerarea
\clearpage

\problem{5.5}{已知解析函数 $f(z)=u(x,y)+\mathrm{i}v(x,y)$ 的实部 $u(x,y)$ 或虚部 $v(x,y)$，求此解析函数。
\begin{enumerate}
  \item $v(x,y)=xy$ \qquad
  \item $v(x,y)=\dfrac{y}{x^2+y^2}$
  \item $u(x,y)=\mathrm{e}^{-y}\sin x$ \qquad
  \item $u(x,y)=\cos x\cosh y$
\end{enumerate}}
\answerarea

\problem{5.6}{在平面极坐标系（$z=\rho\mathrm{e}^{\mathrm{i}\varphi}$）中，已知解析函数 $f(z)=u(\rho,\varphi)+\mathrm{i}v(\rho,\varphi)$ 的实部 $u(\rho,\varphi)=\varphi$，且 $f(1)=0$，求此解析函数。}
\answerarea
\clearpage

\problem{5.7}{在直角坐标系中，已知解析函数 $f(z)=u(x,y)+\mathrm{i}v(x,y)$ 的实部和虚部之和为 $u(x,y)+v(x,y)=-(x+y)(x^2-4xy+y^2)$，求此解析函数。}
\answerarea

\problem{5.8}{有两条平行相距为 $2a$ 且均匀带电的细金属线，每单位长度所带的电荷量分别是 $+q$ 和 $-q$，求与此两金属线垂直的平面静电场的复势、电场线和等势线。}
\answerarea
\clearpage

\sproblem{5.9}{试分析以下几种常见的解析函数变换。
\begin{enumerate}
  \item 线性变换：$W=f(z)=az+b,\quad a,b\in R,\text{ 且 }a\neq 0$；
  \item 幂函数变换：$W=f(z)=z^n,\quad n\in\mathbf{Z}$；
  \item 根式变换：$W=f(z)=z^{\frac{1}{n}},\quad n\in\mathbf{Z}$；
  \item 对数变换：$W=f(z)=\ln z$；
  \item 指数变换：$W=f(z)=\mathrm{e}^z$；
  \item 根式变换：$W=f(z)=\dfrac{az+b}{cz+d},\quad a,b,c,d\in R\text{ 且 }ad-bc\neq 0$。
\end{enumerate}}
\answerarea
\clearpage

\chapter{复变函数积分}

\problem{6.1}{请计算 $f(z)=z^*$ 的回路积分
\[
  \oint_{|z|=1} z^*\,\mathrm{d}z
\]
所得的结果说明了什么问题。}
\answerarea

\problem{6.2}{请计算 $f(z)=\dfrac{1}{z^2-1}$ 的如下回路积分：
\begin{enumerate}
  \item $\displaystyle\oint_{C_1} \frac{1}{z^2-1}\,\mathrm{d}z$ \quad $C_1$ 为 $|z|<1$ 的以原点为圆心的圆周；
  \item $\displaystyle\oint_{C_2} \frac{1}{z^2-1}\,\mathrm{d}z$ \quad $C_2$ 为 $|z|>1$ 的以原点为圆心的圆周；
\end{enumerate}
这两个积分的结果说明了什么问题。}
\answerarea
\clearpage

\chapter{复变函数的级数展开}

\problem{7.1}{求下列幂级数的收敛半径。
\begin{enumerate}
  \item $\displaystyle\sum_{k=1}^{\infty}\left(\frac{z}{k}\right)^k$ \qquad
  \item $\displaystyle\sum_{k=1}^{\infty} k!\left(\frac{z}{k}\right)^k$ \qquad
  \item $\displaystyle\sum_{k=1}^{\infty} k^k z^k$
  \item $\displaystyle\sum_{k=1}^{\infty} k^{\ln k} z^k$ \qquad
  \item $\displaystyle\sum_{k=1}^{\infty} \frac{k^{\ln k}}{k!} z^k$
\end{enumerate}}
\answerarea

\problem{7.2}{已知幂级数 $\displaystyle\sum_{k=0}^{\infty} a_k z^k$ 和 $\displaystyle\sum_{k=0}^{\infty} b_k z^k$ 的收敛半径分别为 $R_a$ 和 $R_b$，求下列幂级数的收敛半径。
\begin{enumerate}
  \item $\displaystyle\sum_{k=0}^{\infty}(a_k-b_k)z^k$ \qquad
  \item $\displaystyle\sum_{k=0}^{\infty}\frac{a_k}{b_k}z^k,\quad (b_k\neq 0)$
\end{enumerate}}
\answerarea
\clearpage

\problem{7.3}{在 $z=0$ 点，将下列函数展成 Taylor 级数，并求出其收敛半径。
\begin{enumerate}
  \item $f(z)=\dfrac{\mathrm{e}^z}{1-z}$ \qquad
  \item $f(z)=\mathrm{e}^{\frac{1}{z-1}}$ \qquad
  \item $f(z)=\arctan z$
  \item $f(z)=\displaystyle\int_0^z \mathrm{e}^{z^2}\,\mathrm{d}z$ \qquad
  \item $f(z)=\displaystyle\int_0^z \frac{\sin z}{z}\,\mathrm{d}z$
\end{enumerate}}
\answerarea

\problem{7.4}{求函数 $f(z)=\sin^3 z$ 在 $z=0$ 点的展开式。}
\answerarea
\clearpage

\problem{7.5}{把函数 $f(z)=\dfrac{1}{z-b}$ 展成幂级数 $\displaystyle\sum_{k=0}^{\infty} c_k (z-a)^k$ 的形式，并给出可以作此展开的条件，其中 $a,b$ 为不相等的复常数。}
\answerarea

\problem{7.6}{请指出下列函数的奇点，并说明属于什么类型的奇点，对于极点请说明极点的阶数。
\begin{enumerate}
  \item $\dfrac{z-1}{z(z^2+1)^2}$ \qquad
  \item $\dfrac{1}{(z^2+\mathrm{i})^2}$ \qquad
  \item $\dfrac{\mathrm{e}^z-1}{z^m}\quad m\text{ 为正整数}$
  \item $\dfrac{\mathrm{e}^z}{\mathrm{e}^z-1}$ \qquad
  \item $\dfrac{1-\cos z}{z^2}$ \qquad
  \item $\dfrac{1-\cos z}{z^3}$
  \item $\dfrac{1}{\sin z}$ \qquad
  \item $\dfrac{1}{\sin z+\cos z}$ \qquad
  \item $\dfrac{\tan(z-1)}{z-1}$
  \item $\cos\dfrac{1}{z+\mathrm{i}}$
\end{enumerate}}
\answerarea
\clearpage

\problem{7.7}{把函数 $f(z)=\dfrac{1}{z^2-3z+2}$ 在下列区域中展成 Laurent 级数
\begin{enumerate}
  \item $0<|z|<1$ \qquad
  \item $1<|z|<2$ \qquad
  \item $2<|z|<\infty$
  \item $0<|z-1|<1$ \qquad
  \item $1<|z-1|<\infty$ \qquad
  \item $0<|z-2|<1$
  \item $1<|z-2|<\infty$
\end{enumerate}}
\answerarea

\problem{7.8}{请将函数 $f(z)=\dfrac{1}{z(z-\mathrm{i})}$ 分别在 $z=1,z=0,z=\mathrm{i}$ 点进行 Laurent 展开，把对每个展开点所有可能的环形区域上的 Laurent 展开式写出。}
\answerarea
\clearpage

\problem{7.9}{请将函数 $f(z)=\dfrac{1}{z^4(z^2-1)^2}$ 在 $z=0$ 点的所有环形区域上展成 Laurent 级数。}
\answerarea

\problem{7.10}{请将函数 $f(z)=\dfrac{1}{(z-1)(z-2)^2}$ 分别在 $z=0,z=1$ 和 $z=2$ 点的所有环形区域上展成 Laurent 级数。}
\answerarea
\clearpage

\problem{7.11}{请将函数 $f(z)=\mathrm{e}^{\frac{1}{1-z}}$ 在 $|z|<1$ 和 $|z|>1$ 的区域中展成级数。}
\answerarea

\problem{7.12}{请将函数 $f(z)=\sin\dfrac{1}{z}$ 在 $|z|>1$ 的区域中展成级数。}
\answerarea
\clearpage

\chapter{留数定理及其在实积分中的应用}

\problem{8.1}{求下列函数在孤立奇点的留数（其中 $n,m$ 为自然数），并求出函数孤立奇点的留数之和。
\begin{enumerate}
  \item $f(z)=\dfrac{1}{z(z-1)^2}$ \qquad
  \item $f(z)=z^n\sin\dfrac{1}{z}$
  \item $f(z)=\dfrac{\mathrm{e}^z-1}{z^5}$ \qquad
  \item $f(z)=\dfrac{\mathrm{e}^{z^2}}{(z-1)^2}$
  \item $f(z)=\dfrac{\mathrm{e}^z}{(z-1)^2}$ \qquad
  \item $f(z)=\mathrm{e}^{\frac{1}{1-z}}$
  \item $f(z)=\dfrac{z^2}{\cos z-1}$ \qquad
  \item $f(z)=z^3\cos\dfrac{1}{z-2}$
  \item $f(z)=\dfrac{1}{(z-\alpha)(z-\beta)^m}$ \qquad
  \item $f(z)=\dfrac{1}{1+z^{2m}}$
  \item $f(z)=\dfrac{z^{2m}}{1+z^m}$ \qquad
  \item $f(z)=\dfrac{\tg z}{1-\mathrm{e}^z}$
\end{enumerate}}
\answerarea

\problem{8.2}{计算下列回路积分
\begin{enumerate}
  \item $\displaystyle\oint_C \frac{1}{(z^2+1)(z-1)^2}\,\mathrm{d}z \qquad C:\ |z-1|=1$
  \item $\displaystyle\oint_C \frac{1}{(z^2+1)(z-1)^2}\,\mathrm{d}z \qquad C:\ |z-\mathrm{i}|=1$
  \item $\displaystyle\oint_C \frac{1}{(z^2+1)(z-1)^2}\,\mathrm{d}z \qquad C:\ |z|=2$
  \item $\displaystyle\oint_{|z|=2} \frac{1}{(z-3)(z^4-1)}\,\mathrm{d}z$ \qquad
  \item $\displaystyle\oint_{|z|=3} \cot^3 z\,\mathrm{d}z$
  \item $\displaystyle\oint_{|z|=1} \frac{z\sin z}{(1-\mathrm{e}^z)^3}\,\mathrm{d}z$ \qquad
  \item $\displaystyle\oint_{|z|=1} z^n\mathrm{e}^{\frac{1}{z}}\,\mathrm{d}z$
\end{enumerate}}
\answerarea
\clearpage

\problem{8.3}{计算下列积分
\begin{enumerate}
  \item $\displaystyle\int_0^{2\pi} \frac{1}{2+\cos\theta}\,\mathrm{d}\theta$ \qquad
  \item $\displaystyle\int_0^{2\pi} \cos^{2n}\theta\,\mathrm{d}\theta$
  \item $\displaystyle\int_0^{2\pi} \frac{1}{1+\sin^2\theta}\,\mathrm{d}\theta$ \qquad
  \item $\displaystyle\int_0^{2\pi} \frac{1}{1-2p\cos\theta+p^2}\,\mathrm{d}\theta,\quad p>0$
\end{enumerate}}
\answerarea

\problem{8.4}{计算下列积分
\begin{enumerate}
  \item $\displaystyle\int_0^\infty \frac{1}{x^4+1}\,\mathrm{d}x$ \qquad
  \item $\displaystyle\int_0^\infty \frac{x^2}{(x^2+a^2)^2}\,\mathrm{d}x$
  \item $\displaystyle\int_0^\infty \mathrm{e}^{-ax^2}\cos bx\,\mathrm{d}x,\quad a>0,\,b\in R$
  \item $\displaystyle\int_0^\infty \frac{(\cos 2x)^2}{(x^2+1)^2(x^4+4)}\,\mathrm{d}x$ \qquad
  \item $\displaystyle\int_{-\infty}^\infty \frac{x^2}{(x^2+1)^2(x^2+2x+2)}\,\mathrm{d}x$
  \item $\displaystyle\int_{-\infty}^\infty \frac{\mathrm{e}^{ax}}{1+\mathrm{e}^x}\,\mathrm{d}x,\quad 0<a<1$
\end{enumerate}}
\answerarea
\clearpage

\problem{8.5}{计算下列积分
\begin{enumerate}
  \item $\displaystyle\int_0^\infty \frac{\cos mx}{1+x^4}\,\mathrm{d}x,\quad m>0$ \qquad
  \item $\displaystyle\int_{-\infty}^\infty \frac{x\sin x}{x^2+a^2}\,\mathrm{d}x,\quad a>0$
  \item $\displaystyle\int_0^\infty \frac{\sin^2 x}{x^2}\,\mathrm{d}x$ \qquad
  \item $\displaystyle\int_{-\infty}^\infty \frac{x\cos x}{x^2-2x+10}\,\mathrm{d}x$
  \item $\displaystyle\int_0^\infty \frac{\sin mx}{x(x^2+a^2)}\,\mathrm{d}x,\quad a>0,\,m\text{ 为正整数}$ \qquad
  \item $\displaystyle\int_0^\infty \frac{1}{1+x^3}\,\mathrm{d}x$
\end{enumerate}}
\answerarea

\problem{8.6}{计算下列积分
\begin{enumerate}
  \item $\displaystyle P\int_{-\infty}^\infty \frac{1}{x^3-1}\,\mathrm{d}x$ \qquad
  \item $\displaystyle P\int_{-\infty}^\infty \frac{\mathrm{e}^{\mathrm{i}kx}-1}{x^2}\,\mathrm{d}x,\quad k>0$
  \item $\displaystyle P\int_{-\infty}^\infty \frac{1}{x(x-1)(x-2)}\,\mathrm{d}x$ \qquad
  \item $\displaystyle P\int_0^\infty \frac{\sin(x+a)\sin(x-a)}{x^2-a^2}\,\mathrm{d}x$
\end{enumerate}}
\answerarea
\clearpage

\problem{8.7}{在下列情况下计算积分 $\displaystyle\int_{-\infty}^\infty \frac{1}{x^2-a^2}\,\mathrm{d}x,\quad a>0$
\begin{enumerate}
  \item 让 $a\to a+\mathrm{i}\varepsilon,\quad \varepsilon\to 0,\,\varepsilon>0$。
  \item 让 $a\to a-\mathrm{i}\varepsilon,\quad \varepsilon\to 0,\,\varepsilon>0$。
  \item 直接计算该积分的 Cauchy 主值积分。
\end{enumerate}}
\answerarea
\clearpage


\part{积分变换与 $\delta$ 函数}
\chapter{Fourier（傅里叶）变换}

\problem{9.1}{如果 $f(x)$ 是实函数，即 $f(x)=f^*(x)$，当 $f(x)$ 展成如下形式的 Fourier 级数时
\[
  f(x) = \sum_{n=-\infty}^{\infty} c_n \mathrm{e}^{\mathrm{i}nx}
\]
系数 $c_n$ 要满足什么条件，请证明。}
\answerarea

\problem{9.2}{请证明：
\[
  \sum_{n=1}^{\infty} (-1)^{n+1}\frac{\sin nx}{n} = \frac{x}{2}
\]}
\answerarea
\clearpage

\problem{9.3}{将如下函数展成 Fourier 级数
\begin{enumerate}
  \item $f(x)=x,\quad x\in[0,2l]$ 展成以 $2l$ 为周期的 Fourier 级数。
  \item $f(x)=x,\quad x\in[0,l]$ 展成以 $2l$ 为周期的 Fourier 级数。
  \item $f(x)=|x|,\quad x\in[-l,l]$ 展成以 $2l$ 为周期的 Fourier 级数。
\end{enumerate}}
\answerarea

\problem{9.4}{电压模式为 $E(t)=E_0\sin\omega t$（$E_0$ 为常数）的交流电。
\begin{enumerate}
  \item 经过半波整流后，其电压的模式为：
  \[
    E(t) = \begin{cases}
      0 & t\in\left[-\dfrac{\pi}{\omega},\,0\right] \\[8pt]
      E_0\sin\omega t & t\in\left[0,\,\dfrac{\pi}{\omega}\right]
    \end{cases}
  \]
  \item 经过全波整流后，其电压模式为
  \[
    E(t) = E_0|\sin\omega t|
  \]
\end{enumerate}
试把以上这两种整流后的电压模式，展成 Fourier 级数，并对这两种展开的结果，从直流成分、基波成分及高次谐波部分等几方面进行比较。}
\answerarea
\clearpage

\problem{9.5}{把振幅按双曲线衰减的振动
\[
  f(t) = \frac{1}{t}\sin\omega t
\]
展成 Fourier 积分，并作出原函数与像函数的简单示意图。}
\answerarea

\problem{9.6}{把定义在 $(0,\infty)$ 上的函数 $f(x)=\mathrm{e}^{-ax},\ (a>0)$，在边界为 $f(0)=0$ 的条件下，展成 Fourier 积分。}
\answerarea
\clearpage

\problem{9.7}{把定义在 $(0,\infty)$ 上的函数 $f(x)=1-H(x-a),\ (a>0)$，在边界为 $f'(0)=0$ 的条件下，展成 Fourier 积分。}
\answerarea

\problem{9.8}{试用 Fourier 变换，求解有阻尼的谐振子方程
\[
  \ddot{x}(t) + \beta\dot{x}(t) + \omega_0^2 x(t) = f(t)
\]}
\answerarea
\clearpage

\chapter{Laplace（拉普拉斯）变换}

\problem{10.1}{求下列函数的 Laplace 变换的像函数。
\begin{enumerate}
  \item $\sinh\omega t$ \qquad
  \item $\cosh\omega t$
  \item $\cos(\omega t+\alpha)$ \qquad
  \item $\cos^2\omega t$
  \item $\mathrm{e}^{-\alpha t}\sin\omega t,\quad \alpha>0$ \qquad
  \item $\dfrac{\cos at-\cos bt}{b^2-a^2},\quad a^2\neq b^2$
\end{enumerate}}
\answerarea

\problem{10.2}{若一函数 $f(t)$，其 Laplace 的像函数为 $F(p)$，请证明：
\[
  \lim_{t\to 0^+} f(t) = \lim_{p\to\infty} pF(p)
\]}
\answerarea
\clearpage

\problem{10.3}{已知下列 Laplace 变换的像函数 $F(p)$，求原函数 $f(t)$：
\begin{enumerate}
  \item $F(p)=\dfrac{1}{(p+a)(p+b)},\quad a\neq b$ \qquad
  \item $F(p)=\dfrac{3p}{p^2-1}$
  \item $F(p)=\dfrac{E}{Lp^2+Rp+\dfrac{1}{C}},\quad E,L,R,C\text{ 都为大于零的实常数}$
  \item $F(p)=\dfrac{1}{(p^2+a^2)(p^2+b^2)},\quad a^2\neq b^2$
  \item $F(p)=\dfrac{\mathrm{e}^{-\alpha p}}{p^2},\quad \alpha>0$ \qquad
  \item $F(p)=\dfrac{1}{p}\cdot\dfrac{\mathrm{e}^{-\alpha p}}{1-\mathrm{e}^{-\alpha p}},\quad \alpha>0$
\end{enumerate}}
\answerarea

\problem{10.4}{已知误差函数的定义为：
\[
  \erf(x) = \frac{2}{\sqrt{\pi}}\int_0^x \mathrm{e}^{-\tau^2}\,\mathrm{d}\tau
\]
试求 $\erf(\sqrt{t})$ 的像函数。}
\answerarea
\clearpage

\problem{10.5}{用 Laplace 变换求下列积分：
\begin{enumerate}
  \item $f(t)=\displaystyle\int_0^\infty \frac{\sin tx}{x}\,\mathrm{d}x$ \qquad
  \item $f(t)=\displaystyle\int_0^\infty \frac{\sin tx}{x(x^2+1)}\,\mathrm{d}x$
  \item $\displaystyle\int_0^\infty \mathrm{e}^{-\sqrt{x}}\,\mathrm{d}x$
\end{enumerate}}
\answerarea

\problem{10.6}{用 Laplace 变换求解下列常微分方程，其中 $y=y(t),\,z=z(t)$
\begin{enumerate}
  \item $\left\{\begin{aligned} &y''-4y=0 \\ &y(0)=-1,\ y'(0)=1 \end{aligned}\right.$ \qquad
  \item $\left\{\begin{aligned} &y'''-y''+4y'-4y=t \\ &y(0)=0,\ y'(0)=0,\ y''(0)=1 \end{aligned}\right.$
  \item $\left\{\begin{aligned} &y'''+3y''+3y'+y=6\mathrm{e}^{-t} \\ &y(0)=0,\ y'(0)=0,\ y''(0)=0 \end{aligned}\right.$ \qquad
  \item $\left\{\begin{aligned} &y'+y+2z'+2z=\mathrm{e}^{-t} \\ &3y'-y+4z'+z=0 \\ &y(0)=-1 \\ &z(0)=0 \end{aligned}\right.$
\end{enumerate}}
\answerarea
\clearpage

\problem{10.7}{用 Laplace 变换求解 $T(t)$ 的方程
\begin{enumerate}
  \item $\left\{\begin{aligned} &T''+\frac{\pi^2 a^2}{l^2}T=A\sin\omega t \\ &T(0)=0,\ T'(0)=0 \end{aligned}\right.$ \\
  其中 $a,A,\omega$ 均为正的常数。
  \item $\left\{\begin{aligned} &T'+\omega^2 a^2 T=g(t) \\ &T(0)=0 \end{aligned}\right.$ \\
  其中 $g(t)$ 为已知函数，$\omega,a$ 为正的常数。
\end{enumerate}}
\answerarea
\clearpage

\chapter{$\delta$ 函数}

\problem{11.1}{若函数 $\delta_\varepsilon(x)$ 的定义为
\[
  \delta_\varepsilon = \begin{cases}
    0 & x>\varepsilon,\ x<0 \\[4pt]
    \dfrac{1}{\varepsilon} & 0\leqslant x<\varepsilon
  \end{cases}
\]
请验证 $\lim_{\varepsilon\to 0}\delta_\varepsilon(x)=\delta(x)$}
\answerarea

\problem{11.2}{请验证
\[
  \lim_{n\to\infty}\sqrt{\frac{n}{\pi}}\mathrm{e}^{-nx^2} = \delta(x)
\]}
\answerarea
\clearpage

\problem{11.3}{请验证
\[
  \lim_{n\to\infty}\frac{1}{\pi}\int_0^n \cos kx\,\mathrm{d}k = \lim_{n\to\infty}\frac{\sin nx}{\pi x} = \delta(x)
\]}
\answerarea

\problem{11.4}{设函数 $f(t)=\cos xt$ 的 Laplace 变换为 $F(p)$，即
\[
  F(p) \doteqdot \cos xt
\]
请验证
\[
  \lim_{p\to 0}\frac{1}{\pi}F(p) = \delta(x)
\]}
\answerarea
\clearpage

\problem{11.5}{对 Heaviside 函数
\[
  H(t-\tau) = \begin{cases}
    1 & t\geqslant\tau \\
    0 & t<\tau
  \end{cases}
\]
请验证
\[
  H'(t-\tau) = \delta(t-\tau)
\]
并以此求 $\delta(t-\tau)$ 的 Laplace 变换的像函数。}
\answerarea

\problem{11.6}{已知某一频谱为：
\[
  f(\omega) = \frac{2A\omega_0}{\pi(\omega^2-\omega_0^2)}\sin\!\left(\frac{\omega}{\omega_0}2\pi N\right)
\]
请验证
\[
  \lim_{N\to\infty} f(\omega) = A\delta(\omega-\omega_0) - A\delta(\omega+\omega_0)
\]
其中 $A,\omega_0$ 为正的常数，$N$ 为正整数。}
\answerarea
\clearpage

% \input{chapters/chapter12.tex} % *第十二章 小波变换初步（无习题）

\part{数学物理方程}
\setcounter{chapter}{12}
\chapter{波动方程、输运方程、Poisson（泊松）方程及其定解问题}

\problem{13.1}{一均匀材质做成的细长弹性圆杆，长为 $l$，截面积为 $S$，静止放置在水平 $x$ 轴上，在 $t=0$ 时刻，给杆一个沿杆的纵向扰动，使杆产生了纵向的微小位移，试推导杆的纵向振动方程。

如果在 $t=0$ 时刻，该扰动使杆的纵振动各点离开平衡位置的分布为 $u(x,0)=A\sin\dfrac{\pi x}{l}$，$A$ 为常数，且 $u_t(x,0)=0$，分别对如下四种边界状况，写出相应的杆的振动定解问题。
\begin{enumerate}
  \item 杆的两个端点 $x=0$ 和 $x=l$ 点刚性固定；
  \item $x=0$ 点固定，$x=l$ 点为自由端；
  \item 在 $x=0$ 点按 $A\sin\omega t$ 运动，在 $x=l$ 点按 $A\cos\omega t$ 运动，$A,\omega$ 均为常数；
  \item 两个端点为弹性固定，端点受到了与位移成正比（比例系数为 $k$）但与位移方向相反的纵向作用力。
\end{enumerate}}
\answerarea

\problem{13.2}{一根长为 $l$ 的匀质弦，放置于介质中，当弦振动时，受到介质的阻尼力，请对阻尼力建立一个合理的模型，并推导弦在此介质中有阻尼的弦的振动方程。

请你给出恰当的初始条件，写出相应的四种不同的弦（任给出四种不同的边界条件）的阻尼振动的定解问题。}
\answerarea
\clearpage

\problem{13.3}{用匀质材料做成的细圆锥杆（指圆锥角很小，即锥体的长度很长而底部的圆形截面半径很小），试推导此圆锥杆的纵振动方程。

把所得的方程与 13.1 题的细长圆杆的纵振动方程进行比较，说明两个方程之间的差异。}
\answerarea

\problem{13.4}{一电阻率为 $\gamma$，截面均为 $S$ 的匀质导线，通有电流密度为 $j$ 的均匀分布的直流电，设该材质的导热系数为 $\kappa$，比热容为 $c$，密度为 $\rho$，设导线的侧面均为绝热的，试推导该导线上的热传导方程。}
\answerarea
\clearpage

\problem{13.5}{设一匀质金属球，其半径为 $a$，以球心为坐标原点，请写出球坐标下的热传导方程（设导热系数为 $\kappa$）。

若把球的表面涂黑，置于阳光下，设球面附近阳光的热流强度为 $M$（设阳光为平行光，在垂直于阳光的单位面积上，在单位时间内流经这一单位面积上的太阳光的热量为 $M$），而球外环境的温度为 $0\,^\circ\mathrm{C}$，金属球表面同时按牛顿冷却定律散热。设在 $t=0$ 时刻，金属球的温度为 $T_0$，试给出该定解问题。}
\answerarea

\problem{13.6}{在铀块中，除了中子的扩散运动（扩散系数为 $D$）外，还存在中子的吸收和增殖反应过程。设在单位时间内，单位体积中吸收和增殖的中子数正比于该时刻该处中子的浓度 $u(\boldsymbol{r},t)$，即 $\beta u(\boldsymbol{r},t)$，$\beta$ 为比例常数（或称为净增殖系数）。若我们选取的铀块是匀质的球状铀块，铀块的半径为 $a$。试在球坐标下导出中子浓度 $u(\boldsymbol{r},t)$ 扩散所满足的方程。

若球状铀块的表面边界满足如下几种不同的条件，请写出相应条件下的定解问题。（初始条件请做个合理的假设。）
\begin{enumerate}
  \item 铀块表面的中子浓度为零；
  \item 没有中子从铀块表面逸出（或是从铀块表面逸出的中子数目与回到铀块表面的中子数目相等）；
  \item 铀块表面的中子浓度不随时间变化。
\end{enumerate}}
\answerarea
\clearpage

\problem{13.7}{写出静电场中电介质表面静电场的衔接条件。}
\answerarea

\problem{13.8}{一根导热杆由两段 $l_1$ 与 $l_2$ 相衔接构成，两端的导热系数、比热和密度分别为 $\kappa_1,c_1,\rho_1$ 和 $\kappa_2,c_2,\rho_2$。设杆的初始温度为 $u_0$，保持杆的两个端点的温度为 $0\,^\circ\mathrm{C}$。试写出该导热杆的热传导的定解问题（假定导热杆的侧面均为绝热的）。}
\answerarea
\clearpage

\problem{13.9}{在均匀介质区域 $\Omega$ 内，对在稳定的热流分布 $f(\boldsymbol{r})$ 产生的温度场 $u(\boldsymbol{r})$ 满足 Poisson 方程 $\nabla^2 u(\boldsymbol{r})=-f(\boldsymbol{r})$，其中 $f(\boldsymbol{r})=\dfrac{F(\boldsymbol{r})}{\kappa}$，$\kappa$ 为导热系数，$F(\boldsymbol{r})$ 是热流密度。
\begin{enumerate}
  \item 试说明第一，二，三类边界条件在此问题中的物理意义；
  \item 试说明第二类边界条件存在的必要条件。
\end{enumerate}}
\answerarea

\problem{13.10}{在真空区域 $\Omega$ 内，有稳定的电荷分布 $\rho(\boldsymbol{r})$，则在 $\Omega$ 中存在稳定的电势分布，电势 $u(\boldsymbol{r})$ 满足 Poisson 方程
\[
  \nabla^2 u(\boldsymbol{r}) = -\frac{1}{\varepsilon_0}\rho(\boldsymbol{r})
\]
其中 $\varepsilon_0$ 为真空介电常数。若 $\Omega$ 的边界 $\partial\Omega$ 是一导体壳，当
\begin{enumerate}
  \item 导体壳接地；
  \item 导体壳不接地。
\end{enumerate}
分别对这两种情况，说明第一类和第二类边界条件的物理意义。}
\answerarea
\clearpage

\chapter{分离变量法}

\problem{14.1}{长为 $l$ 的弦，两端固定，弦中的张力为 $T$，在距弦的一端 $x=0$ 的 $l_0$ 处（$l_0<l$），以力 $F_0$ 把弦拉开，使 $l_0$ 点离开平衡位置有一小位移，并使弦处于稳定状态。在 $t=0$ 时刻突然撤销这一大小为 $F_0$ 的力，求解弦的振动。}
\answerarea

\problem{14.2}{长为 $l$ 的弦，两端固定，在弦的 $x_0$ 处（$0<x_0<l$）施予谐变外力 $\rho f(t)=\rho A\sin\omega t$，（$\rho$ 为弦的线密度，$A$ 为常数），求解弦的振动。

（提示：外力可表示为 $\rho f(x,t)=\rho A\delta(x-x_0)\sin\omega t$。）}
\answerarea
\clearpage

\problem{14.3}{求解习题 13.1 中的（2）、（3）两种情况下杆的振动。}
\answerarea

\problem{14.4}{导热系数为 $\kappa$ 的均匀物质构成的立方体，$0\leqslant x\leqslant l_1,\ 0\leqslant y\leqslant l_2,\ 0\leqslant z\leqslant l_3$，其初始温度为 $T_0$，如果其表面温度保持在 $0\,^\circ\mathrm{C}$，求此立方体温度的变化。并写出此立方体中心点，即 $\left(\dfrac{l_1}{2},\dfrac{l_2}{2},\dfrac{l_3}{2}\right)$ 处温度随时间的变化关系。又问当时间 $t\to\infty$ 时，该立方体的温度表达式，并从物理上解释其合理性。

思考：如果边界条件改为该立方体置于 $0\,^\circ\mathrm{C}$ 的环境中，其边界以牛顿冷却定律与环境进行热交换，其解的形式又是如何变化的，并与上面的结果进行比较，分析这两种边界条件的物理图像上的差异。}
\answerarea
\clearpage

\problem{14.5}{侧面绝热的均匀细杆，$0\leqslant x\leqslant l$，其导热系数为 $\kappa$，初始温度为 $T_0$，其两端保持不变的温度：$u(0,t)=T_1,\,u(l,t)=T_2$，$T_1,T_2$ 均为常数。求杆上的温度分布随时间的变化关系。当 $t\to\infty$ 时，给出杆上的温度分布的表达式，并解释这一结果在物理上的合理性及其满足的物理定律。}
\answerarea

\problem{14.6}{铀块中中子的扩散、吸收的增殖反应如习题 13.4 所述，试求一维细铀杆的临界长度。临界长度是指当铀杆超过这一长度时，杆中的中子浓度随时间而增长致使铀块产生爆炸。这里忽略杆的过中心轴剖面的对角线长度与杆的长度之差。}
\answerarea
\clearpage

\problem{14.7}{一厚度为 $l$ 的单晶硅片置于充满杂质的环境中，杂质从硅片的两表面向里扩散，假设硅片两表面杂质浓度保持恒定的 $N_0$，杂质在硅片中的扩散系数为 $D$，求解此扩散问题。并说明扩散的时间与扩散系数 $D$ 及厚度 $l$ 的关系。当 $t$ 较大时，化简所得的结果。}
\answerarea

\problem{14.8}{在矩形区域 $0<x<a,\ 0<y<b$ 上求解 Laplace 方程
\[
  \nabla^2 u(x,y) = 0
\]
使其满足边界条件：
\[
  \left\{\begin{aligned}
    u(0,y) &= Ay(b-y), & u(a,y) &= 0 \\
    u(x,0) &= B\sin\frac{\pi x}{a}, & u(x,b) &= 0
  \end{aligned}\right.
\]
其中 $A,B$ 为常数。}
\answerarea
\clearpage

\problem{14.9}{在量子物理学中，空间中自由粒子可视为 de Broglie（德布罗意）波，被称为物质波，用 $\psi(\boldsymbol{r},t)$ 来描述，它满足如下 Schrödinger（薛定谔）方程（量子波动方程）
\[
  \mathrm{i}\hbar\frac{\partial}{\partial t}\psi(\boldsymbol{r},t) = -\frac{\hbar^2}{2m}\nabla^2\psi(\boldsymbol{r},t)
\]
其中 $\hbar$ 为 Planck（普朗克）常数，$m$ 为粒子质量。我们把粒子看作在一个三边长为 $L$ 的正立方体 $\left(-\dfrac{L}{2}\leqslant x\leqslant\dfrac{L}{2},\ -\dfrac{L}{2}\leqslant y\leqslant\dfrac{L}{2},\ -\dfrac{L}{2}\leqslant z\leqslant\dfrac{L}{2}\right)$ 中的三维波动，且这一波动满足周期性边界条件，试求解这一波动。}
\answerarea
\clearpage

\chapter{曲线坐标系下方程的分离变量}

\problem{15.1}{试用平面极坐标 $(\rho,\varphi)$，把二维波动方程分离变量。}
\answerarea

\problem{15.2}{在量子物理学中氢原子满足的定态 Schrödinger 方程
\[
  -\frac{\hbar^2}{2m}\nabla^2 u(\boldsymbol{r}) - \frac{e^2}{r}u(\boldsymbol{r}) = E u(\boldsymbol{r})
\]
其中 $\hbar,m,e,E$ 都为常数，试用球坐标将此方程分离变量。}
\answerarea
\clearpage

\problem{15.3}{在 $x=0$ 的邻域内求下列方程的两个独立的级数解，并对这两个方程的级数解进行比较
\begin{enumerate}
  \item $y''-xy'=0$；
  \item $y''-x^2y'=0$。
\end{enumerate}}
\answerarea

\problem{15.4}{在 $x=0$ 的邻域内求解方程
\[
  (x^2-1)y'' + xy' - y = 0
\]
的两个独立的级数解。}
\answerarea
\clearpage

\problem{15.5}{在 $x=0$ 的邻域内求 Chebyshev（切比雪夫）方程
\[
  (1-x^2)y'' - xy' + n^2 y = 0\quad (n\text{ 为正整数})
\]
的级数解。}
\answerarea

\problem{15.6}{在 $x=0$ 的邻域内，用级数展开法求解 Euler-Cauchy 方程
\[
  x^2 y'' + 2x y' = 0
\]}
\answerarea
\clearpage

\problem{15.7}{在 $x=0$ 的邻域内求方程
\[
  xy'' - xy' + y = 0
\]
的两个独立的级数解。}
\answerarea

\problem{15.8}{在 $z=0$ 的邻域内求解下列方程的两个独立级数解，并对这两个方程的级数解进行比较
\begin{enumerate}
  \item $\dfrac{\mathrm{d}^2 u(z)}{\mathrm{d}z^2} + \dfrac{2}{z}\dfrac{\mathrm{d}u(z)}{\mathrm{d}z} + m^2 u(z) = 0$；
  \item $\dfrac{\mathrm{d}^2 u(z)}{\mathrm{d}z^2} + \dfrac{1}{z}\dfrac{\mathrm{d}u(z)}{\mathrm{d}z} - m^2 u(z) = 0$。
\end{enumerate}}
\answerarea
\clearpage

\problem{15.9}{对二阶齐次常微分方程
\[
  y'' + p(x)y' + q(x)y = 0
\]
设 $y_1$ 和 $y_2$ 是它的两个线性独立的解，并可定义 Wronske（朗斯基）行列式 $W(x)$
\[
  W(x) = y_1 y_2' - y_1' y_2 = \begin{vmatrix} y_1 & y_2 \\ y_1' & y_2' \end{vmatrix}
\]
\begin{enumerate}
  \item 试证明 $W(x)$ 满足
  \[
    W'(x) = -p(x)W(x)
  \]
  \item $W(x)$ 可写为
  \[
    W(x) = y_1^2 \frac{\mathrm{d}}{\mathrm{d}x}\left(\frac{y_2}{y_1}\right)
  \]
  试证明：两个解 $y_1(x)$ 和 $y_2(x)$ 有如下关系
  \[
    y_2(x) = C y_1(x) \int_b^x \frac{\mathrm{e}^{-\int_a^{x'} p(x'')\,\mathrm{d}x''}}{[y_1(x')]^2}\,\mathrm{d}x'
  \]
  其中 $a,b$ 为任意常数，$C$ 为积分常数，它与 $a,b$ 的选取有关。上式表明，对于本题所给定的二阶齐次常微分方程的形式，当已知一个解 $y_1(x)$ 时，可通过上式求出另一个线性独立的解 $y_2(x)$。
\end{enumerate}}
\answerarea

\problem{15.10}{利用 15.9 题的结果，已知下列方程的一个解，求出另一个线性独立的解，
\begin{enumerate}
  \item 已知方程
  \[
    \frac{\mathrm{d}^2 y}{\mathrm{d}x^2} + y = 0
  \]
  的一个解 $y_1=\sin x$，求另一个解 $y_2$。
  \item 已知方程
  \[
    \frac{\mathrm{d}^2 R(r)}{\mathrm{d}r^2} + \frac{1}{r}\frac{\mathrm{d}R(r)}{\mathrm{d}r} - \frac{m^2}{r^2}R(r) = 0
  \]
  的一个解 $R_1(r)=r^m$，求另一个解 $R_2(r)$。
\end{enumerate}}
\answerarea
\clearpage

\problem{15.11}{考虑 Sturm-Liouville 本征值的问题：对于方程
\[
  \frac{1}{w(x)}\frac{\mathrm{d}}{\mathrm{d}x}\left[w(x)a(x)\frac{\mathrm{d}u(x)}{\mathrm{d}x}\right] + [c(x)-\lambda]u(x) = 0,\quad x\in[a,b]
\]
其中函数 $w(x),a(x)$ 和 $c(x)$ 满足如下条件：
\begin{gather*}
  a(x) = a_0 + a_1 x + a_2 x^2 \\
  \frac{1}{w(x)}\frac{\mathrm{d}}{\mathrm{d}x}[w(x)a(x)] = b_0 + b_1 x \\
  c(x) = c_0 \\
  [w(x)a(x)]_{x=a} = [w(x)a(x)]_{x=b} = 0
\end{gather*}
且 $w(x)a(x)$ 在区间 $(a,b)$ 内不为零。$a_0,a_1,a_2,b_0,b_1,c_0$ 均为实数。

请证明：对于给定的本征值 $\{\lambda_n\}$
\[
  \lambda_n = n(n-1)a_2 + nb_1 + c_0,\quad n=0,1,2,3,\cdots
\]
\begin{enumerate}
  \item 总存在非平凡的多项式解
  \[
    u_n(x) = \sum_{k=0}^n \alpha_k^{(n)} x^k
  \]
  其中 $\alpha_k^{(n)}$ 为 $n$ 次多项式的展开系数。
  \item 这些多项式具有正交性
  \[
    \int_a^b u_m^*(x)u_n(x)w(x)\,\mathrm{d}x = 0,\quad m\neq n
  \]
\end{enumerate}}
\answerarea

\problem{15.12}{在 $x=0$ 点邻域内，求 Laguerre（拉盖尔）方程
\[
  xy'' + (1-x)y' + \lambda y = 0
\]
的级数解。当 $\lambda$ 取什么值时，可使级数退化为多项式？这些多项式乘以适当的常数使最高幂次项成为 $(-x)^n$ 的形式，此时的多项式称为 Laguerre 多项式，记为 $L_n(x)$，写出 $L_n(x)$ 的前三项。}
\answerarea
\clearpage

\problem{15.13}{在 $x=0$ 点邻域内，求解 Gauss 方程（超几何级数方程），
\[
  x(x-1)y'' + [(1+\alpha+\beta)x-\gamma]y' + \alpha\beta y = 0
\]
的级数解。}
\answerarea
\clearpage

\chapter{球函数}

\problem{16.1}{试证明 Legendre 多项式的生成函数
\[
  f(r) = \frac{1}{\sqrt{1+r^2-2rx}}
\]
在 $r=0$ 点的 Taylor 展开式为
\[
  f(r) = \sum_{l=0}^{\infty} \mathrm{P}_l(x)r^l
\]
其中 $r=|\boldsymbol{r}|<1,\,x=\cos\theta,\,|x|\leqslant 1$。}
\answerarea

\problem{16.2}{试证明 Legendre 多项式的递推公式：
\begin{enumerate}
  \item $(2l+1)\mathrm{P}_l(x)=\mathrm{P}_{l+1}'(x)-\mathrm{P}_{l-1}'(x)$
  \item $(l+1)\mathrm{P}_l(x)=\mathrm{P}_{l+1}'(x)-x\mathrm{P}_l'(x)$
  \item $l\mathrm{P}_l(x)=x\mathrm{P}_l'(x)-\mathrm{P}_{l-1}'(x)$
  \item $(x^2-1)\mathrm{P}_l'(x)=lx\mathrm{P}_l(x)-l\mathrm{P}_{l-1}(x)$
\end{enumerate}}
\answerarea
\clearpage

\problem{16.3}{试证明：
\begin{enumerate}
  \item $\mathrm{P}_l(1)=1$
  \item $\mathrm{P}_l(-1)=(-1)^l$
  \item $\mathrm{P}_l(0)=\begin{cases} 0 & l=2n+1,\ n=0,1,2,\cdots \\[4pt] (-1)^n\dfrac{(2n-1)!!}{2n!!} & l=2n,\ n=0,1,2,\cdots \end{cases}$
  \item $\mathrm{P}_l'(1)=\dfrac{l(l+1)}{2}$
  \item $\mathrm{P}_l'(-1)=(-1)^{l-1}\dfrac{l(l+1)}{2}$
  \item $\mathrm{P}_l'(0)=\begin{cases} (-1)^n\dfrac{(2n+1)!}{(2n)!!}, & l=2n+1,\ n=0,1,2,\cdots \\[6pt] 0, & l=2n,\ n=0,1,2,\cdots \end{cases}$
\end{enumerate}}
\answerarea

\problem{16.4}{计算下列积分：
\begin{enumerate}
  \item $\displaystyle\int_{-1}^{1} \mathrm{P}_l(x)\,\mathrm{d}x$ \qquad
  \item $\displaystyle\int_{-1}^{1} x^2\mathrm{P}_l(x)\mathrm{P}_{l+2}(x)\,\mathrm{d}x$
  \item $\displaystyle\int_{-1}^{1} (1-x^2)[\mathrm{P}_l'(x)]^2\,\mathrm{d}x$ \qquad
  \item $\displaystyle\int_{-1}^{1} x\mathrm{P}_l(x)\mathrm{P}_{l+1}(x)\,\mathrm{d}x$
  \item $\displaystyle\int_{0}^{1} \mathrm{P}_l(x)\,\mathrm{d}x$ \qquad
  \item $\displaystyle\int_{0}^{1} \mathrm{P}_k(x)\mathrm{P}_l(x)\,\mathrm{d}x$
  \item $\displaystyle\int_{-1}^{1} x^k\mathrm{P}_l(x)\,\mathrm{d}x$，分别讨论 $k\geqslant l$ 和 $k<l$ 的情况。
\end{enumerate}}
\answerarea
\clearpage

\problem{16.5}{在以 Legendre 多项式为基 $\{\mathrm{P}_l(x),\,l=0,1,2,\cdots\}$ 的函数空间中，将下列函数展开。
\begin{enumerate}
  \item $f(x)=|x|$， \qquad
  \item $f(x)=1+x+x^2+x^3$
\end{enumerate}}
\answerarea

\problem{16.6}{在点电荷 $4\pi\varepsilon_0 q$ 的电场中放置一半径为 $a$ 的导体球壳，球心距点电荷的距离为 $b$，且 $b>a$，求这体系所形成的电场。}
\answerarea
\clearpage

\problem{16.7}{半径为 $a$ 的匀质半球，其半球面的温度保持恒定为 $T_0$，其底面温度保持零度，求球内温度的稳定分布。}
\answerarea

\problem{16.8}{求以下两种均匀带电体在空间形成的稳定电场分布。
\begin{enumerate}
  \item 均匀带电的薄圆盘，圆盘的半径为 $a$，总电荷量为 $4\pi\varepsilon_0 q$。
  \item 均匀带电的细圆环，环的半径为 $a$，总电荷量为 $4\pi\varepsilon_0 q$。
\end{enumerate}}
\answerarea
\clearpage

\problem{16.9}{求解半径为 $a$ 的空心球壳内的定解问题
\[
  \left\{\begin{aligned}
    &\nabla^2 u = 0,\quad r<a \\
    &u\big|_{r=a} = \cos^2\theta,\quad 0\leqslant\theta\leqslant\pi
  \end{aligned}\right.
\]}
\answerarea

\problem{16.10}{试证明
\[
  \mathrm{P}_l^{-m}(x) = (-1)^m \frac{(l-m)!}{(l+m)!}\mathrm{P}_l^m(x),\quad m>0
\]}
\answerarea
\clearpage

\problem{16.11}{试证明
\[
  \mathrm{P}_l^m(0) = \begin{cases}
    (-1)^{\frac{1}{2}(l-m)}\dfrac{(l+m-1)!!}{(l-m)!!}, & \text{当 }l+m\text{ 为偶数} \\[8pt]
    0, & \text{当 }l-m\text{ 为奇数}
  \end{cases}
\]}
\answerarea

\problem{16.12}{试证明
\[
  \mathrm{P}_l^l(\cos\theta) = (2l-1)!!\sin^l\theta,\quad l=0,1,2,\cdots
\]}
\answerarea
\clearpage

\problem{16.13}{在以球函数为基 $\left\{\mathrm{Y}_{l,m}(\theta,\varphi),\,\begin{aligned} &l=0,1,2,\cdots \\ &m=0,\pm 1,\cdots,\pm l \end{aligned}\right\}$ 的函数空间中，将下列函数展开
\begin{enumerate}
  \item $\sin^2\theta\cos^2\varphi$ \qquad
  \item $(1+3\cos\theta)\sin\theta\cos\varphi$
  \item $\dfrac{1}{r^2}(x^2+2z^2+3xy+4xy)$
\end{enumerate}
其中 $r^2=x^2+y^2+z^2,\,0\leqslant\theta\leqslant\pi,\,0\leqslant\varphi\leqslant 2\pi$。}
\answerarea

\problem{16.14}{在半径为 $a$ 的球壳内外，分别求解下面定解问题的球外区域和球内区域的解
\[
  \left\{\begin{aligned}
    &\nabla^2 u = 0 \\
    &u\big|_{r=a} = f(\theta,\varphi) \\
    &u\big|_{r=0}<\infty,\quad u\big|_{r\to\infty}<\infty
  \end{aligned}\right.
\]
并分别对以下两种 $f(\theta,\varphi)$ 的形式
\begin{enumerate}
  \item $f(\theta,\varphi)=\cos\theta$，
  \item $f(\theta,\varphi)=\mathrm{P}_l^1(\cos\theta)\cos\varphi$，给出解的具体形式。
\end{enumerate}}
\answerarea
\clearpage

\chapter{柱函数}

\problem{17.1}{试证明
\begin{gather*}
  \cos x = \mathrm{J}_0(x) + 2\sum_{m=1}^{\infty}(-1)^m\mathrm{J}_{2m}(x) \\
  \sin x = 2\sum_{m=0}^{\infty}(-1)^m\mathrm{J}_{2m+1}(x)
\end{gather*}
（提示：对整数阶 Bessel 函数的生成函数 $\mathrm{e}^{\frac{x}{2}\left(z-\frac{1}{z}\right)}$ 取 $z=\mathrm{i}$ 进行展开运算。）}
\answerarea

\problem{17.2}{试证明
\begin{gather*}
  \frac{1}{\pi}\int_0^\pi \cos(x\sin\theta)\cos m\theta\,\mathrm{d}\theta = \begin{cases} \mathrm{J}_m(x), & m\text{ 为偶数} \\ 0, & m\text{ 为奇数} \end{cases} \\
  \frac{1}{\pi}\int_0^\pi \sin(x\sin\theta)\sin m\theta\,\mathrm{d}\theta = \begin{cases} 0, & m\text{ 为偶数} \\ \mathrm{J}_m(x), & m\text{ 为奇数} \end{cases}
\end{gather*}
并进一步证明
\[
  \mathrm{J}_m(x) = \frac{1}{\pi}\int_0^\pi \cos(m\theta-x\sin\theta)\,\mathrm{d}\theta,\quad m=0,1,2,\cdots
\]
（提示：对整数阶 Bessel 函数的生成函数 $\mathrm{e}^{\frac{x}{2}\left(z-\frac{1}{z}\right)}$ 取 $z=\mathrm{e}^{\mathrm{i}\theta}$ 进行展开运算。）}
\answerarea
\clearpage

\problem{17.3}{试证明
\[
  1 = [\mathrm{J}_0(x)]^2 + 2\sum_{m=1}^{\infty}[\mathrm{J}_m(x)]^2
\]
因此有结论 $|\mathrm{J}_0(x)|^2\leqslant 1,\quad |\mathrm{J}_m(x)|\leqslant\dfrac{1}{\sqrt{2}}$

（提示：对整数阶 Bessel 函数的生成函数 $\mathrm{e}^{\frac{x}{2}\left(z-\frac{1}{z}\right)}$ 和 $\mathrm{e}^{\frac{x}{2}\left[(-z)-\frac{1}{(-z)}\right]}$ 相乘可证。）}
\answerarea

\problem{17.4}{用整数阶 Bessel 函数的生成函数证明
\[
  \mathrm{J}_m(x) = (-1)^m\mathrm{J}_m(-x)
\]
这说明整数阶 Bessel 函数的对称性，当 $m$ 为奇数时为奇对称，当 $m$ 为偶数时为偶对称。注意：此关系式由 $\mathrm{J}_m(x)$ 的级数表达式可直接得到。}
\answerarea
\clearpage

\problem{17.5}{计算下列积分
\begin{enumerate}
  \item $\displaystyle\int_0^b \mathrm{J}_0(x)\,\mathrm{d}x$ \qquad
  \item $\displaystyle\int_0^b \mathrm{J}_1(x)\,\mathrm{d}x$
  \item $\displaystyle\int_0^b x\mathrm{J}_0(x)\,\mathrm{d}x$ \qquad
  \item $\displaystyle\int_0^b x^3\mathrm{J}_0(x)\,\mathrm{d}x$
  \item $\displaystyle\int_0^b x^4\mathrm{J}_1(x)\,\mathrm{d}x$ \qquad
  \item $\displaystyle\int_0^b \mathrm{J}_3(x)\,\mathrm{d}x$
\end{enumerate}}
\answerarea

\problem{17.6}{在 $\rho\in(0,b)$ 的区间上，在第一类齐次边界条件下，在由零阶 Bessel 函数所构成的本征函数空间中（该空间的基向量为 $\left\{\mathrm{J}_0\!\left[\dfrac{x_n^{(0)}}{b}\rho\right],\,n=1,2,\cdots\right\}$，$x_n^{(0)}$ 为 $\mathrm{J}_0$ 的第 $n$ 个零点），把函数 $f(\rho)=1$ 展开。}
\answerarea
\clearpage

\problem{17.7}{在高为 $h$，半径为 $b$ 的柱体内，取柱坐标的原点在下底的圆心，对于边界条件为
\[
  \left\{\begin{aligned}
    &Z(z)\big|_{z=0} = Z(z)\big|_{z=h} = 0 \\
    &R(\rho)\big|_{\rho=b} = 0
  \end{aligned}\right.
\]
求解稳定场 $u(\rho,\varphi,z)$ 满足方程 $\nabla^2 u(\rho,\varphi,z)=f(\rho,\varphi,z)$ 的场的分布，并对解得结果说明其物理图像。

求解输运方程
\[
  \frac{\partial u}{\partial t} - a^2\nabla^2 u = f(\rho,\varphi,z)
\]
并对解得结果说明其物理图像。}
\answerarea

\problem{17.8}{边缘固定半径为 $R$ 的圆形膜，初始形状为旋转抛物面
\[
  u\big|_{t=0} = \left(1-\frac{\rho^2}{R^2}\right)H\quad (H\text{ 为常数})
\]
且初速度为零，求此膜的振动。}
\answerarea
\clearpage

\problem{17.9}{高为 $h$，半径为 $b$，导热系数为 $\kappa$ 的圆柱体，其下底面保持零度，上底面的温度为 $\rho$ 的函数 $f(\rho)$，其侧面在零度的环境中自由冷却（满足牛顿冷却定律），冷却系数为 $H$，求柱内温度场的稳定分布。}
\answerarea

\problem{17.10}{半径为 $b$ 介电常数为 $\varepsilon$ 的长直均质介质圆柱，放入均匀电场 $\boldsymbol{E}_0$ 中，圆柱的轴线与 $\boldsymbol{E}_0$ 垂直，求解柱体内外的电场。}
\answerarea
\clearpage

\problem{17.11}{由圆弧 $\rho=a$ 和 $\rho=b\,(b>a)$，和矢径角度 $\varphi=0$ 与 $\varphi=\alpha$ 围成的扇形区域内（如习题 17.11 图所示），求解 Laplace 方程
\[
  \nabla^2 u = 0
\]
使场 $u(\rho,\varphi)$ 的分布满足边界条件
\[
  \left\{\begin{aligned}
    &u\big|_{\rho=a} = 0,\quad u\big|_{\rho=b} = u_0 \\
    &u\big|_{\varphi=0} = u\big|_{\varphi=\alpha} = 0
  \end{aligned}\right.
\]
\fig[0.42]{17_11}
分别以角度 $\varphi$ 的方程的本征值问题和以径向 $\rho$ 的方程的本征值问题进行求解，比较这两种解法的差异。}
\answerarea

\problem{17.12}{半径为 $r_0$ 的均质球，球面始终保持零度，初始时刻球内温度分布为：
\begin{enumerate}
  \item $u\big|_{t=0}=f(r)$
  \item $u\big|_{t=0}=f(r)\cos\theta$
\end{enumerate}
分别求解这两种初始条件下球内温度分布的变化，当 $t\to\infty$ 时，这两种情况解得的结果进行比较，说明其物理图像。}
\answerarea
\clearpage

\chapter{Green（格林/绿）函数法}

\problem{18.1}{两个相距为 $a$ 均接地的理想导体平面，两平行导体平面间放置一电荷量为 $q$ 的点电荷，电荷距离下平面为 $b,\,b<a$，求导体间的电势分布。（提示：用镜像法求解。）}
\answerarea

\problem{18.2}{一半径为 $a$ 的导体球壳接地，在球外距球心为 $b$ 处（$b>a$）放置一点电荷 $q$，试用镜像法求电场的分布，并求出导体球壳表面感应电荷的分布，进而证明此电荷分布等价于球内的镜像电荷。}
\answerarea
\clearpage

\problem{18.3}{用 Green 函数法求解第一类边值问题的球内 Laplace 方程的解。
\[
  \left\{\begin{aligned}
    &\nabla^2 u = 0,\quad r<a \\
    &u\big|_{r=a} = \cos\theta
  \end{aligned}\right.
\]}
\answerarea

\problem{18.4}{半径为 $b$ 的无穷长均匀圆柱体，初始温度为 $U_0$，初始时刻外界空间的温度为零，用 Green 函数法求解 $t$ 时刻柱外空间的温度 $u(\rho,t),\,\rho>b$。}
\answerarea
\clearpage

\problem{18.5}{侧面为绝热的无穷长细杆，初始温度为零，在 $t=0$ 时刻在 $x=x_0$ 点有一瞬时热源 $Q$，试构造该一维问题的 Green 函数。}
\answerarea
\clearpage

\chapter{其他求解方法及方程}

\problem{19.1}{分别用 Fourier 变换法和 Laplace 变换法求解下面一维无界空间中的输运问题，
\[
  \left\{\begin{aligned}
    &u_t - a^2 u_{xx} = 0,\qquad -\infty<x<\infty \\
    &u\big|_{t=0} = \varphi(x),
  \end{aligned}\right.
\]}
\answerarea

\problem{19.2}{用 Fourier 变换法求解一维无界弦的受迫振动问题
\[
  \left\{\begin{aligned}
    &u_{tt} - a^2 u_{xx} = f(x,t) \\
    &u\big|_{t=0} = \varphi(x) \\
    &u_t\big|_{t=0} = \nu(x)
  \end{aligned}\right.
\]}
\answerarea
\clearpage

\problem{19.3}{设有两根匀质同质的半无界杆，侧面均绝热，温度分别为 $0$ 和 $u_0$，在 $t=0$ 时刻，将两杆的端点相接，求 $t>0$ 时杆中各点的温度分布。}
\answerarea

\problem{19.4}{无界静止弦在 $t=0$ 时刻，在 $x=x_0$ 处受到冲量为 $I$ 的冲击作用，试求解弦的振动。

（提示：由 $t=0^-$ 到 $t=0$ 时刻动量的变化等于冲量，可得 $u_t\big|_{t=0}=\dfrac{I}{\rho}\delta(x-x_0)$，并可由 D'Alembert 公式求解。）}
\answerarea
\clearpage

\problem{19.5}{场量 $u(x,t)$ 满足下列方程，可设方程有行波解 $u(\xi)$，其中 $\xi=x-at$，试把下列方程化为行波变量 $\xi$ 的常微分方程
\begin{enumerate}
  \item $\dfrac{\partial u}{\partial t} + u\dfrac{\partial u}{\partial x} + \beta\dfrac{\partial^3 u}{\partial x^3} = 0$
  \item $\dfrac{\partial u}{\partial t} + u\dfrac{\partial u}{\partial x} + \alpha\dfrac{\partial^2 u}{\partial x^2} + \beta\dfrac{\partial^3 u}{\partial x^3} + \gamma\dfrac{\partial^4 u}{\partial x^4} = 0$
  \item $\dfrac{\partial u}{\partial t} - \alpha\dfrac{\partial^2 u}{\partial x^2} - \gamma u(1-u) = 0$
  \item $\dfrac{\partial^2 u}{\partial t^2} - a^2\dfrac{\partial^2 u}{\partial x^2} + \alpha u^2 = 0$
  \item $\dfrac{\partial^2 u}{\partial t^2} - a^2\dfrac{\partial^2 u}{\partial x^2} - \alpha\dfrac{\partial^4 u}{\partial x^4} - \beta\dfrac{\partial^2 u^2}{\partial x^2} = 0$
\end{enumerate}}
\answerarea

\problem{19.6}{
\begin{enumerate}
  \item 请把 $x^{2k},\,k=0,1,2,\cdots,$ 展成 Hermite 多项式的偶次项级数
  \[
    x^{2k} = \frac{(2k)!}{2^{2k}}\sum_{n=0}^{k}\frac{\mathrm{H}_{2n}(x)}{(2n)!(k-n)!}
  \]
  \item 请把 $x^{2k+1},\,k=0,1,2,\cdots,$ 展成 Hermite 多项式的奇次项级数
  \[
    x^{2k+1} = \frac{(2k+1)!}{2^{2k+1}}\sum_{n=0}^{k}\frac{\mathrm{H}_{2n+1}(x)}{(2n+1)!(k-n)!}
  \]
\end{enumerate}}
\answerarea
\clearpage

\problem{19.7}{请证明：
\begin{enumerate}
  \item $\displaystyle\int_{-\infty}^{\infty} \mathrm{H}_n(x)\mathrm{e}^{-\frac{1}{2}x^2}\,\mathrm{d}x = \begin{cases}
    \dfrac{2\pi n!}{\left(\dfrac{n}{2}\right)!}, & n\text{ 为偶数} \\[12pt]
    0, & n\text{ 为奇数}
  \end{cases}$
  \item $\displaystyle\int_{-\infty}^{\infty} x\mathrm{H}_n(x)\mathrm{e}^{-\frac{1}{2}x^2}\,\mathrm{d}x = \begin{cases}
    0, & n\text{ 为偶数} \\[6pt]
    \dfrac{2\pi (n+1)!}{\left(\dfrac{n+1}{2}\right)!}, & n\text{ 为奇数}
  \end{cases}$
\end{enumerate}}
\answerarea

\problem{19.8}{请证明：
\[
  \int_{-\infty}^{\infty} x^2\mathrm{e}^{-x^2}\mathrm{H}_n(x)\mathrm{H}_n(x)\,\mathrm{d}x = \sqrt{\pi}\,2^n n!\left(n+\frac{1}{2}\right)
\]}
\answerarea
\clearpage

% \input{chapters/chapter20.tex} % *第二十章 非线性数学物理方程初步

\part{变分法初步}
% \input{chapters/chapter21.tex} % 第二十一章 泛函的变分（无习题）
\setcounter{chapter}{21}
\chapter{变分原理}

\problem{22.1}{对于 $y'=\dfrac{\mathrm{d}y}{\mathrm{d}x}\neq 0$，请证明以下两种形式的 Euler-Lagrange 方程
\[
  \frac{\partial f}{\partial x} - \frac{\mathrm{d}}{\mathrm{d}x}\frac{\partial f}{\partial y'} = 0
\]
和
\[
  \frac{\partial f}{\partial y} - \frac{\mathrm{d}}{\mathrm{d}x}\left(f - y'\frac{\partial f}{\partial y'}\right) = 0
\]
是等价的。（Euler-Lagrange 方程具有多种形式。）}
\answerarea

\problem{22.2}{对常参数 $\alpha$ 的泛函
\[
  I(\alpha) = \int_{x_0}^{x_1} f[y(x,\alpha),\,y_x(x,\alpha),\,x]\,\mathrm{d}x
\]
把被积函数 $f$ 在 $\alpha=0$ 点展成参数 $\alpha$ 的 Taylor 级数，由泛函极值的条件 $\dfrac{\partial I(\alpha)}{\partial\alpha}=0$，请导出 Euler-Lagrange 方程。}
\answerarea
\clearpage

\problem{22.3}{求下列泛函取极值的 Euler-Lagrange 方程，并对方程进行求解。
\begin{enumerate}
  \item $I=\displaystyle\int_{x_0}^{x_1} \frac{x}{x+y'}\,\mathrm{d}x$ \qquad
  \item $I=\displaystyle\int_{x_0}^{x_1} \sqrt{1+x}\sqrt{1+{y'}^2}\,\mathrm{d}x$
\end{enumerate}}
\answerarea

\problem{22.4}{求下列不含 $x$ 的泛函取极值时的 Euler-Lagrange 方程，并对方程进行求解。
\begin{enumerate}
  \item $I=\displaystyle\int_{x_0}^{x_1} \sqrt{1+y^2{y'}^2}\,\mathrm{d}x$ \qquad
  \item $I=\displaystyle\int_{x_0}^{x_1} (y^2+{y'}^2)\,\mathrm{d}x$
\end{enumerate}}
\answerarea
\clearpage

\problem{22.5}{请写出粒子作自由落体下落的 Hamilton 量，并给出这一系统的正则方程，说明这一正则方程与牛顿力学中自由落体运动方程的等价性。}
\answerarea

\problem{22.6}{请写出线性谐振子的 Hamilton 量，并求出这一系统的 Hamilton 正则方程。进一步说明这一正则方程与牛顿力学中得到的线性谐振子的运动方程的等价性。}
\answerarea
\clearpage

\problem{22.7}{对于广义动能
\[
  T = \alpha_{ij}\dot{q}_i\dot{q}_j,\qquad \alpha_{ij}=\alpha_{ji},
\]
请证明：
\[
  \frac{\partial T}{\partial\dot{q}_i}\dot{q}_i = 2T
\]}
\answerarea

\problem{22.8}{试构造一个简单的物理系统的作用量，当该系统对时间 $t$ 进行无穷小平移变换时，若系统的作用量保持不变，请证明这一系统的能量是个守恒量。}
\answerarea
\clearpage

\end{document}
